\songtitle{Två år på Söder}

\tag*{2026}\tag{punk-rock}\tag{swedish-lyrics}\tag{two-years-in-sodermalm}
\begin{notes}
Punk-rock re-genre of \emph{Två år på Söder}, the Swedish-language adaptation of \emph{Two Years in Södermalm} (page \pageref{sec:two-years-in-sodermalm}); same lyrics. Retrieved from Gabriel's Suno page (V6-MINI, 2 October 2026).
\end{notes}
\begin{style}
Punk-rock, early-1990s groove: roomy snare ambience, driving mid-tempo backbeat, warm and direct male lead vocal, overdriven power-chord guitar with bass-led verses and a full-band chorus.
\end{style}
\begin{lyrics}
\songpart{Vers 1}
Jag såg dig på avstånd i snart två år,\\
kring Medborgarplatsen, där tunnelbanan går,\\
på väg till vuxengymnasiet vid Gullmarsplan,\\
ibland på andra håll i stan.

\songpart{Vers 2}
Jag fick aldrig veta vem du var, aldrig ditt namn,\\
bara din stil, samma upproriska punk varje säsong,\\
ändå bara en fast punkt i stadens sammanhang.

\annotation{Refräng}\\
Två år på Söder, jag gick och jag gick,\\
en fast punkt i människornas landskap:\\
en rödhårig punkare, en rödtott med blick,\\
som jag nästan blev kär i.

\songpart{Vers 3}
Du höll mig förtrollad i snart två år,\\
två år såg jag dig i kvarteren där vi går,\\
i ett kvarter för stort för att du skulle se mig,\\
länets skeva mitt, där ingen ser varandra.

\annotation{Refräng}\\
Två år på Söder, jag gick och jag gick,\\
en fast punkt i människornas landskap:\\
en rödhårig punkare, en rödtott med blick,\\
som jag nästan blev kär i.

\annotation{Brygga}\\
Samma jacka, samma strumpor, året runt,\\
samma störande punkestetik,\\
en mask över en vacker kvinna därinne\\
som jag inte kunde låta bli att längta efter.

\annotation{Sista refräng}\\
Två år på Söder, jag gick och jag gick,\\
en fast punkt i människornas landskap:\\
en rödhårig punkare,\\
en vacker kvinna under,\\
ett långt längtande, en lång dröm,\\
en flicka som jag nästan blev kär i.
\end{lyrics}

